\documentclass{revtex4-2}
\usepackage[utf8]{inputenc}
\usepackage[T1]{fontenc}
\usepackage{geometry}
\usepackage{graphicx}
\usepackage{float}
\usepackage{hyperref}
\usepackage[english,indonesian]{babel}

\title{Optimasi Pemrosesan Template LaTeX Menggunakan Pendekatan Berbasis AI}
\author{Academic Support Agent}
\date{\today}

\begin{document}
\maketitle

\begin{abstract}

\end{abstract}

\chapter{PENDAHULUAN}


\chapter{TINJAUAN PUSTAKA}


% Original content (commented out):
% 
%%%%%%%%%%%%%%%%%%%%%%%%%%%%%%%%%%%%%%%%%%%%%%%%%%%%%%%%
%  Plantilla para reportes de Física Experimental
%  Intermedia I de la Licenciatura en Física en la
%  Universidad Autónoma Metropolitana - Iztapalapa
%
%  Preparado por Roberto Olayo-Valles
%  E-mail: rolv@xanum.uam.mx
%  Febrero 2019
%%%%%%%%%%%%%%%%%%%%%%%%%%%%%%%%%%%%%%%%%%%%%%%%%%%%%%%%

% Esta plantilla se basa en el paquete REVTeX de la
% American Physical Society (APS).

% Primero se declara el tipo de documento:

\documentclass[aps,prb,twocolumn, amsmath, amssymb]{revtex4-2}

% Para esta plantilla utilizaremos la plantilla de la
% revista Physical Review B de la APS. Por eso
% utilizamos las opciones aps y prb. La opción
% twocolumn formateará el documento en dos columnas.
% amsmath y amssymb nos permite usar ciertos símbolos
% matemáticos.

% A continuación declaramos los paquetes que necesitamos

\usepackage{graphicx}  % Para insertar figuras

\usepackage[utf8]{inputenc}  % Para que las letras en
                             % español aparezcan
                             % correctamente

\usepackage[spanish,mexico]{babel} % Para que el texto
                                   % insertado aparezca
                                   % en español
                                   
\usepackage{dcolumn}	% Para poder alinear columnas
                        % en tablas con el punto decimal
\usepackage{bm}  % Para poner símbolos matemáticos en
       	         % negritas

% Las siguientes instrucciones modifican el texto
% que aparece antes de la fecha en la plantilla.
\makeatletter
\def\Dated@name{Fecha de preparación: }
\makeatother

% Comienzo del documento
\begin{document}

% Título
\title{Escritura de Reportes de Física Experimental Intermedia I en \LaTeX}

% Información del autor (si hay más de un autor,
% repertir los comandos con su información)
\author{Roberto Olayo-Valles}
\email{rolv@xanum.uam.mx}
\affiliation{Departamento de Física, Universidad Autónoma Metrpolitana\textendash Iztapalapa}

% Fecha
\date{\today}

% Resumen
\begin{abstract}
	Este documento tiene tres propósitos: servir como plantilla para la escritura de reportes de Física Experimental Intermedia I, señalar el contenido que debe tener cada sección del reporte, y mostrar ejemplos del uso de \LaTeX. Esta sección es el resumen, y aunque es la primer sección del documento, generalmente es la última en escribirse. Debes describir de manera resumida la motivación, el procedimiento y el (o los) resultados. Por lo general, no se incluyen ecuaciones ni citas bibliográficas en el resumen.
\end{abstract}

% Instrucción para insertar el encabezado con título,
% autores, fecha y resumen:
\maketitle

% Sección de "Introducción"
\section{\label{sec:intro}Introducción}
Como parte de tu formación en la Licenciatura en Física, debes aprender a comunicar los resultados de tu trabajo. En particular, en este curso pondremos énfasis en la presentación escrita. La importancia de preparar informes o reportes bien escritos no es menor. En palabras de José Miró Julià \cite{Miro}:

\begin{quote}
	Quizá la habilidad más importante que pueda tener una persona es saber escribir. Los empresarios, jefes, directores esperan pedir a sus empleados que estudien un tema y redacten un informe, y poderlo leer y entender con facilidad; es esencial en el trabajo en equipo para poder tener constancia escrita de lo que va haciendo cada miembro del equipo; ayuda a clarificar ideas cuando se trabaja un tema. Es importante saber escribir, pues al final casi lo único que queda de nuestro trabajo es lo que hayamos escrito.
\end{quote}

En este curso aprenderás a preparar reportes en una de las herramientas más utilizadas en el campo de la Física para producir documentos de calidad profesional: \LaTeX (que se pronucia \textit{latec} o a veces \textit{leitec} en inglés, pero nunca \textit{lateks}). En el aula virtual de nuestro curso se encuentran los archivos fuente con que se creó este documento. Es importante que revises el archivo fuente y lo vayas modificando hasta entender cómo utilizar \LaTeX.

Esta plantilla está hecha a partir del paquete REV\TeX\ de la American Physical Society y sigue el formato de la revista \textit{Physical Review B}. De manera que la apariencia de tus reportes será similar a la de artículos científicos publicados en dicha revista. \textbf{\LaTeX\ se encargará de darle el formato adecuado a tu reporte; tú sólo debes preocuparte por el contenido.} Esta es una forma totalmente diferente de escribir un reporte respecto a escribirlo en un procesador de texto como Microsoft Word.

% Subsección
\subsection{\label{sec:latex}Escribiendo en \LaTeX}
La mayor parte de tu texto lo escribirás con el uso normal del teclado con algunas excepciones:

% Lista
\begin{itemize}
	\item Para escribir comillas utiliza dos acentos graves (\texttt{\`}\texttt{\`}) 
	para abrir y dos apostrofes (\texttt{''}) 
	para cerrar. Por ejemplo, para escribir ``texto'', en el código fuente debes escribir \texttt{\`}\texttt{\`}\texttt{texto''}.
	\item Los símbolos \&, \%. \$, \#, \_, \{ y \} 
	son caracteres especiales en \LaTeX. Para escribirlos debes poner una diagonal invertida antes del símbolo. Por ejemplo, \verb|\%| para escribir \%.
	\item Otros símbolos especiales en \LaTeX\ son \textasciitilde, \textasciicircum \space y \textbackslash. Para escribir éstos debes usar \verb|\textasciitilde|, \verb|\textasciicircum| y \verb|\textbackslash|, respectivamente.
\end{itemize}

La forma de la letra puede modificarse utilizando los comandos que se muestran en la Tabla~\ref{tab:formas_texto}. Después del comando coloca el texto al que le quieres aplicar la forma entre llaves \{\}. Por ejemplo, para escribir \textit{Physical Review B}, en el código fuente debes escribir \verb|\textit{Physical Review B}|.

% Tabla
\begin{table}
	\caption{\label{tab:formas_texto}Comandos de forma de texto en \LaTeX}
	\begin{ruledtabular}
		\begin{tabular}{ll}
			\multicolumn{2}{c}{\textbf{Formas de texto normal}}\\
			\textbf{Comando} & \textbf{Forma}\\
			\verb|textit{texto}| & Cursiva\\
			\verb|textbf{texto}| & Negritas\\
			\verb|texttt{texto}| & Máquina (Typewriter)\\
			\verb|emph{texto}| & Énfasis\\
			\verb|underline{texto}| & Subrayado\\
		\end{tabular}
	\end{ruledtabular}
\end{table}

\subsection{\label{sec:refs}Referencias Cruzadas}

En el texto de tu reporte comúnmente harás referencia a alguna sección previa, ecuaciones, tablas o figuras. En \LaTeX estas referencias se hacen con los comandos \verb|\label{referencia}| y \verb|\ref{referencia}|. Por ejemplo, la tabla de ``Comandos de forma de texto en \LaTeX'' contiene el comando \verb|\label{tab:formas_texto}| en su estructura en el archivo fuente. Si se quiere hacer referencia a la tabla, se utilizaría algo como: ``Véase la Tabla\verb|~\ref{tab:formas_texto}|''. El resultado es el siguiente: ``Véase la Tabla~\ref{tab:formas_texto}.

\subsection{\label{sec:cyb}Citas y Bibliografía}

Especialmente importante en la escritura de reportes es citar trabajos previos y dar crédito a quien corresponde. Todos los científicos nos basamos en el trabajo de otros y es importante citar esos trabajos. Para incluir una cita en el texto debes escribir \verb|\cite{id}| donde \verb|id| es un identificador de la referencia. Por ejemplo, en el primer párrafo de la Introducción de este documento cité un documento de José Miró; para ello escribí \verb|\cite{Miro}| en donde quiero que se coloque el número de cita.

Si deseas citar varias referencias en el mismo lugar en el texto, separa sus identificadores con comas: \verb|\cite{id1,id2,id3}|. Cuando compiles el documento, \LaTeX se encargará de generar la numeración adecuadamente.

La bibliografía debe colocarse al final del documento (justo antes de \verb|\end{document}|). Para colocarla, escribimos la instrucción \verb|\bibliography{nombre_archivo_bib}|, donde \verb|nombre_archivo_bib| es el nombre del archivo que contiene los datos de la bibliografía.

El archivo  que contiene los datos de la bibliografía es un archivo de texto con extensión \verb|.bib|. Este archivo de texto contiene los detalles de la bibliografía con los identificadores de cada trabajo. A continuación se presentan algunos ejemplos de referencias de libros, sección de un libro, página web y artículo científico:

\begin{verbatim}
@BOOK{Bire82,
author       = {N. D. Birell and 
                   P. C. W. Davies},
year         = 1982,
title        = {Quantum Fields in Curved Space},
publisher    = {Cambridge University Press}
}

@INBOOK{Beutler1994,
author       = "E. Beutler", 
editor       = "E. Beutler and 
                 M. A. Lichtman and 
                 B. W. Coller and 
                 T. S. Kipps", 
booktitle    = "Williams Hematology", 
chapter      = "7", 
pages        = "654--662",
publisher    = "McGraw-Hill", 
year         = "1994", 
edition      = "5", 
address      = "New York", 
volume       = "2", 
}

@misc{knuthwebsite,
author    = "Donald Knuth",
title     = "Knuth: Computers and Typesetting",
howpublished = {\url{http://www-cs-faculty.stanford.
                     edu/\~{}uno/abcde.html}}
}

@ARTICLE{feyn54,
author       = "R. P. Feynman",
year         = "1954",
journal      = "Phys.\ Rev.",
volume       = "94",
pages        = "262",
}


\end{verbatim}

También puedes revisar el archivo \verb|.bib| de este documento para ver ejemplos.

Al compilar tu archivo para producir el documento, el compilador utilizará un par de paquetes llamados \verb|natbib| y Bib\TeX para producir la bibliografía. Es importante que sepas que existe una alternativa a Bib\TeX llamada \verb|biber|; éste paquete no funciona para compilar la bibliografía de esta plantilla (y en general de documentos que usan el paquete REV\TeX). Por lo tanto, cuando compiles el documento asegúrate de usar Bib\TeX en las opciones del compilador.

\section{\label{sec:exp}Experimental}

En esta sección describes los materiales, aparatos y  procedimientos utilizados en tu experimento. Generalmente, se escribe en tiempo verbal pasado. Es muy útil en esta sección incluir una figura que muestre la disposición del experimento. Para incluir figuras utilizarás el ambiente \texttt{figure}. En este ambiente le asignarás una descripción a la figura y \LaTeX se encargará de asignarle un número. El ambiente se ve de la siguiente manera:

\begin{verbatim}
\begin{figure}
	\includegraphics[width=\linewidth]{fig1.pdf}
	\caption{\label{fig1}Texto de descripción de
	 la figura.}
\end{figure}
\end{verbatim}

Es recomendable que la figura esté en formato pdf, aunque el paquete \texttt{graphicx} funciona también con otros formatos. Si quieres que la figura tenga un ancho menor al ancho de la columna puedes colocar una fracción antes de \verb|\linewitdth|. Por ejemplo, si quieres que la figura ocupe el 70\% del ancho de la columna escribirías \verb|[width=0.7\linewidth]|. Si la figura es muy ancha y prefieres que ocupe todo el ancho de la página utiliza \verb|figure*|. 

A continuación un ejemplo: Conectamos las resistencias y la fuente para formar un circuito como se muestra en la Figura~\ref{fig:circuito}.

\begin{figure}
	\includegraphics[width=0.7\linewidth]{figura_1}
	\caption{\label{fig:circuito}Esquema del circuito armado para el experimento. En la descripción de la figura puedes}
\end{figure}

\section{\label{sec:ryd}Resultados y Discusión}

En esta sección mostrarás tus resultados en tablas y figuras (generalmente gráficas). Ya se mostró en la sección \ref{sec:exp} como insertar figuras. Para insertar una tabla utilizarás el ambiente \verb|table| (de manera similar al caso de figuras, si necesitas poner una tabla que abarque el ancho de las dos columnas utiliza \verb|table*|). Este ambiente lleva más elementos que el de \verb|figure|. A continuación está un ejemplo de los comandos usados para generar una tabla:

\begin{verbatim}
\begin{table}
\caption{\label{tab:param_eh}Parámetros 
    utilizados en el experimento de 
    electrohilado}
\begin{ruledtabular}
\begin{tabular}{cccc}
Concentración & Velocidad de flujo & Voltaje 
	& Distancia\\
(g/ml) & (ml/h) & (kV) & (cm)\\
\colrule
17.0 & 0.3 & 18 & 10.0 \\
18.0 & 0.4 & 21 & 15.0 \\
19.0 & 0.5 & 25 & 20.0 \\
\end{tabular}
\end{ruledtabular}
\end{table}
\end{verbatim}

El resultado se puede ver en la Tabla \ref{tab:param_eh}. Dentro del ambiente \verb|table|, la primer instrucción es \verb|caption| que funciona para ponerle título a la tabla e incluye \verb|ref| para poder hacer referencia a ella en el texto. Después se coloca el ambiente \verb|ruledtabular| que pone las líneas dobles en la parte superior e inferior de la tabla. A continuación sigue el ambiente \verb|tabular| seguido entre corchetes de una serie de letras. El número de letras define el número de columnas que tiene la tabla (cuatro en el ejemplo) y las posibles letras son: \texttt{l}, \texttt{c}, \texttt{r} y \texttt{d} para que el contenido de la columna está alineado a la izquierda, centrado, derecha o con el punto decimal, respectivamente. Finalmente, se encuentra el contenido de la tabla. Para cada fila se escriben los elementos separados por signos \&, y al final de la fila se coloca \textbackslash \textbackslash. Los encabezados de las columnas se separan del resto de la tabla con una línea horizontal; para ello utilizamos \verb|\colrule|.

\begin{table}
	\caption{\label{tab:param_eh}Parámetros utilizados en el experimento de electrohilado}
	\begin{ruledtabular}
		\begin{tabular}{cccc}
			Concentración & Velocidad de flujo & Voltaje & Distancia\\
			(g/ml) & (ml/h) & (kV) & (cm)\\
			\colrule
			17.0 & 0.3 & 18 & 10.0 \\
			18.0 & 0.4 & 21 & 15.0 \\
			19.0 & 0.5 & 25 & 20.0 \\
		\end{tabular}
	\end{ruledtabular}
\end{table}


\subsection{\label{sec:mate}Expresiones Matemáticas y Ecuaciones}
Una de las características más importantes de \LaTeX es que permite escribir expresiones matemáticas de manera correcta y consistente. Si quieres incluir una expresión matemática como parte del texto, debes colocar la expresión entre símbolos \$. Por ejemplo, para escribir $F = ma$, escribes \verb|$F = ma$|. Además de expresiones matemáticas, también se usa esta estrategia para referirte a variables dentro del texto. Por ejemplo: La energía $E$ está relacionada con la masa $m$ por la ecuación $E=mc^2$, donde $c$ es la velocidad de la luz. En ese ejemplo todas las variables y la ecuación fueron colocadas entre símbolos \$.

Para escribir superíndices y subíndices usa \^ y \_. Por ejemplo: \verb|$a_0^n$| resulta en $a_0^n$. Si el subíndice o superíndice tiene más de un caracter, debes encerrar la expresión entre corchetes: para obtener $e^{2\pi}$, escribe \verb|$e^{2\pi}$|. Para expresiones más complejas puedes consultar diversas fuentes en internet \cite{wikimath,ams}

Hay veces que necesitamos incluir texto normal dentro de una expresión matemática, para ello utilizamos la instrucción \verb|\text{texto}|. Por ejemplo, para obtener $y=x \text{ para toda } x_{\text{e-p}}$ escribes \verb|$y=x \text{ para toda } x_{\text{e-p}}$|. Fíjate que \LaTeX se encarga de colocar espacios alrededor de ciertos símbolos como $=$, pero en el texto tú debes colocarlos.

Para ecuaciones numeradas utiliza el ambiente \verb|equation|. Por ejemplo,
\begin{equation}
\frac{\partial \rho}{\partial t} + \nabla \cdot \bm{J} = 0
\label{ec:cons}
\end{equation}
se escribe de la siguiente manera: 
\begin{verbatim}
\begin{equation}
\frac{\partial \rho}{\partial t} + \nabla 
\cdot \bm{J} = 0  
\label{ec:cons}
\end{equation}
\end{verbatim}
La función \verb|\label{}| nuevamente tiene la función de darle una etiqueta a la función para que después podamos hacer referencia a la ecuación.

Si deseas escribir un grupo de ecuaciones relacionadas y quieres que alinearlas por el símbolo de igualdad, puedes usar el ambiente \verb|align|. Por ejemplo, para escribir
\begin{align}
\nabla \cdot \bm{E}& = \frac{\rho}
{\varepsilon_0} \label{ec:Max1}\\
\nabla \times \bm{E}& = -\frac{\partial \bm{B}}
{\partial t} \label{ec:Max2}\\
\nabla \cdot \bm{B}& = 0 \label{ec:Max3}\\
\nabla \times  \bm{B}& = \mu_0 \left( \bm{J} + 
\varepsilon_0 \frac{\partial \bm{E}}{\partial t}
 \right) \label{ec:Max4}
\end{align}
los comandos son los siguientes:
\begin{verbatim}
\begin{align}
\nabla \cdot \bm{E}& = \frac{\rho}
{\varepsilon_0} \label{ec:Max1}\\
\nabla \times \bm{E}& = -\frac{\partial \bm{B}}
{\partial t} \label{ec:Max2}\\
\nabla \cdot \bm{B}& = 0 \label{ec:Max3}\\
\nabla \times  \bm{B}& = \mu_0 \left( \bm{J} + 
\varepsilon_0 \frac{\partial \bm{E}}{\partial t}
\right) \label{ec:Max4}
\end{align}
\end{verbatim}
Si no deseas que las ecuaciones queden alineadas con el símbolo de igualdad, puedes usar el ambiente \verb|gather| en lugar de \verb|align|. Para hacer referencia a alguna ecuación, en lugar de \verb|\ref{}|, utiliza \verb|\eqref{}|. Por ejemplo: La Ecuación~\eqref{ec:cons} es la ley de conservación de carga.

Cabe señalar que en tu reporte no debes poner la derivación detallada de ecuaciones. Especialmente, no debes poner pasos de álgebra.

\subsection{\label{etica}Ética Científica}

Existen varias cuestiones éticas que debes aprender si tu objetivo es ser un científico serio. Aquí expongo dos de estas cuestiones, frabricación de datos y plagio.

Cuando leemos reportes científicos, incluyendo reportes de laboratorio, asumimos que lo descrito en el reporte es una descripción honesta de lo observado por el autor (o autores). Podemos estar en desacuerdo sobre la interpretación de los datos, o la teoría que se propone. Pero confiamos en que si repetimos el experimento en la manera descrita en el reporte, obtendremos los mismos resultados (dentro de la incertidumbre reportada).

Si se encuentra que lo reportado no es reproducible, los autores sufren un golpe importante a su reputación como científicos porque se sospecha fraude. El fraude es una de las faltas más graves que puede cometer un científico porque la comunidad deja de confiar en sus reportes. Cuando esto sucede, se vuelve muy difícil recobrar la confianza y, por lo tanto, la carrera científica de la persona puede quedar arruinada.

Por esta razón, es muy importante que reportes lo que tú hiciste y observaste, sin inventar cosas. Recuerda que en Física nos interesa describir la naturaleza. No se trata de forzar a la naturaleza a comportarse de acuerdo a lo que dicen nuestros libros de texto. Por el contrario, si encontramos que nuestros libros de texto no describen lo que vemos en el experimento, entonces hay que modificar la teoría. En Física y en Ciencia, el experimento (siempre que esté hecho con cuidado) es rey.

La otra falta ética que debes evitar es el plagio. Plagio significa ``copiar en lo sustancial obras ajenas, dándolas como propias''\cite{RAE}. La planeación del curso se especificó la siguiente política sobre copiado:
\begin{quote}
	Copiar el trabajo de otros compañeros, o de fuentes bibliográficas o de Internet, para hacerlo pasar como trabajo propio es una falta grave y no será tolerada. \textbf{Quien incurra en esta falta recibirá calificación de cero en la actividad experimental correspondiente.}
	
	La traducción palabra por palabra de fuentes en otro idioma también se considera copiado y, por lo tanto, ameritará la misma sanción. 
	
	Si tienes duda, consulta conmigo sobre la manera adecuada de citar fuentes.
\end{quote}

\bibliography{reporte_FEI}

\end{document}

\end{document}